\chapter{기계 전체}
% full version: chapters/10_synthesis.tex

\pencilfig{10}{\pencilart{10}}{기계 전체: 데이터 버스, 그 안의 브레이크, 그리고 위에 걸린 제어용 라디오 채널 넷.}

지금까지 뉴런 유형을 하나씩 골라 그것이 기계에 무엇을 더하는지 물었다. 이제 부품을 모아 조립할 차례다.

그렇게 하면 3층짜리 기계가 나온다. 맨 아래에는 \nt{GLUT} 뉴런의 거대한 전화망이 있고, 이 망을 따라 데이터가 흐른다. 그 옆에는 억제성 \nt{GABA} 뉴런이 있다. 이들은 흐름을 제어한다. 무엇을, 언제, 얼마나 크게 통과시킬지 정한다. 그리고 이 모든 것 위에 방송국이 몇 개 있고, 방송국마다 자기 다이얼이 있다. \nt{DOPA}는 연결의 어떤 변화를 남길지 알린다. \nt{SERO}는 전체 수준을 유지하고, 가설에 따르면 인내의 시간 지평도 정한다. \nt{NORA}는 탐색과 결정 사이에서 고른다. \nt{ACET}은 쓰기와 읽기 사이에서 고른다. 실제로 각 방송국은 전선 하나가 아니라 작은 채널 묶음이다. 그래도 860억 개의 뉴런에 채널은 한 줌뿐이다.

\section*{하나의 이야기}

사건 하나가 기계 전체를 거쳐 가는 과정을 따라가 보자. 동물은 오래전부터 알고 있다. 레버 A를 누르면 주스가 나온다. 어느 날 세상이 바뀐다. A는 더 이상 듣지 않고, 주스는 동물이 거의 쓰지 않던 레버 B에서 나온다.

A를 눌렀는데 주스가 오지 않는다. 도파민은 침묵으로 답하고, A를 고른 연결이 약해진다. 침묵이 되풀이되고 오류가 평소보다 커진다. 가설에 따르면 노르아드레날린이 세상이 바뀌었다고 알린다. 대비 다이얼이 내려가고, 선택이 부드러워지며, 잡음이 B를 시도해 보도록 더 자주 떠민다. 아세틸콜린은 네트워크를 쓰기 모드로 바꾼다. B를 눌러 보자 예상 밖의 주스가 나온다. 도파민이 치솟고 B를 고른 연결이 강해진다. 몇 번 시도하고 나면 기대가 새 세상을 따라잡고, 깜짝 놀랄 일이 끝나며, 다이얼들은 제자리로 돌아간다.

어느 다이얼도 다른 다이얼이 무엇을 하는지 모른다는 점에 주목하자. 각 다이얼은 자기 신호에만 반응하고, 동작의 순서는 저절로 짜인다. 각 블록이 입력이 준비되면 동작하는 비동기 회로와 같다. 각 다이얼이 따로따로 이렇게 동작한다는 것은 대부분 측정되었다. 다이얼들이 함께 손발을 맞춰 일한다는 것은 아직 가설이다.

\section*{여러분의 노트북과 다른 점}

노트북에는 공통 클록이 있지만 뇌에는 없다. 노트북의 소자는 믿을 만하지만 시냅스는 신호의 절반 넘게를 잃는다. 노트북의 메모리는 프로세서와 따로 있지만 뇌에서는 메모리가 같은 연결 안에 있다. 노트북은 (학습을 한다면) 오차 역전파로 배우지만, 뇌는 국소 규칙과 “예상보다 좋다”는 브로드캐스트 신호 하나로 배운다. 그리고 가장 중요한 점. 노트북은 프로세서 하나에 수십 와트를 쓰지만 뇌는 모든 것에 20와트를 쓴다.

\section*{우리가 모르는 것}

대뇌 피질이 딸깍의 정확한 타이밍을 얼마나 활용하는지 우리는 모른다. 교사가 모두에게 하나뿐인데 여러 층으로 된 회로의 연결 수십억 개를 어떻게 조정하는지도 모른다. 방송국이 무엇을 전하는지도 완전히 이해하지 못했다. 뇌의 멀리 떨어진 부분들이 공통 클록 없이 어떻게 보조를 맞추는지도 모른다. 그리고 같은 일을 20와트로 해내는 기계를 만들 줄 아는 사람은 아직 아무도 없다.

이 책은 이제 이 핵심을 넘어선다. 기계의 입력과 출력, 기계의 디버깅, 수면, 그리고 칩 위의 살아 있는 뉴런으로 만든 기계로 나아간다.

\fullversion{10}
